\chapter{Background}\label{ch:background}
\label{chap:background}

Since the proposed solution builds upon existing features within modelling and slicing software, it is essential to establish a fundamental understanding of these systems and the file formats connecting them.
The next chapter therefore outlines the key concepts and technologies underlying 3D modelling and slicing, providing the basis for the subsequent analysis and prototype development.

\section{State-of-the-art Workflow}

The following section examines the general workflow described in the previous chapter in greater detail, further using the 3D model illustrated in Figure~\ref{fig:simple-3d-model-box} as an example.

\begin{figure}[h]
    \centering
    \includegraphics[width=\linewidth]{FlowChart_IterativeProcess.png}
    \caption{Iterative workflow in 3D printing highlighting the step to be optimised}
    \label{workflow}
\end{figure}

\begin{enumerate}
    \item 3D Model
    
    After coming up with an idea, the user either creates the 3D model from scratch or downloads one from the various databases available online and optionally modifies it. 
    Creating and modifying a model requires highly domain-specific knowledge and skills.
    Although a wide range of educational resources — including videos, written tutorials, and documentation — are readily available for all levels of expertise, newcomers often require substantial time, practice, and effort before their designs approach the intended outcome.

    \item Print Setup 
    
    As the next step, the user has to choose the printing material, called filament, for different parts, \eg printing the box (\ref{fig:simple-3d-model-box}, A and B) itself in white PLA and the text (C) in a different color PLA, and set at the very least the temperature according to the materials specifications in case they differ from the presets used. 
    Oftentimes further parameters or presets are needed to determine factors like supports\footnote{3D printers build a model by printing one layer on top of another. While overhangs can be printed up to a certain degree and distance, in some cases a supporting structure is needed as basis for upper layers.\cite[p.~152]{gibson_additive_2010}}, infill\footnote{Infill is the structure printed within a hollow 3D model. It can vary regarding the amount, usually indicated in percent, as well as the pattern. Both of these affect the model's resulting stability.\cite[p.~155]{gibson_additive_2010}}, printing speed and heating bed temperature. 

    In this case, the green boxes around the closures (D) need to be marked as modifiers to increase the wall loops in the main model for added stability. 
    Furthermore the box is designed to be printed as two separate bodies, one each for the top and bottom, which means the top including all its text bodies have to be placed on the printing bed accordingly to save on material and avoid support structures. 
    \item Slicing 
    
    The model must then be translated into layer-by-layer instructions for the printer, including the generation of supports and necessary helping structures such as the wiping tower for multi-material printing.
    Although this is a technically complex process, the complexity is handled entirely by the software and therefore does not present significant challenges for the user.
    \item Printing 
    
    After physically setting up the printer, \ie loading the chosen materials and optionally performing calibrations, the model can be printed using the slicers instructions. 
    \item Evaluation
    
    Especially for newcomers and casual makers, the initial print often does not meet expectations.
    The printed object must therefore be evaluated for potential design flaws or printing errors.
    This often leads to revisions to either the 3D model or the printing setup, initiating a new iteration of the workflow. 
\end{enumerate}

Ergo, two parts of the process require a considerable specialised knowledge: the creation of the 3D model and the preparation of the model for printing.
Particularly given the iterative nature of the workflow, several repetitive steps can be identified.

For instance, when multiple materials are used — such as printing text in a different colour — the user has to assign the corresponding material presets to the respective model components in each iteration.
As project complexity increases, this task can become increasingly tedious and prone to error.

A similar challenge arises when model parts are employed as modifiers, for example to strengthen specific areas by increasing wall thickness around a snap-fit closure.
In such cases, the modifiers must also be redefined with every iteration, and unlike material assignments, they cannot always be saved as presets.
This further reinforces the need for systematic documentation and reapplication of settings throughout the process.

\section{3D Modeling Software}

When realising a design from scratch - rather than downloading and modifying an existing model — the user typically employs 3D modelling software.
Such software enables the transformation of an idea or a two-dimensional sketch into a three-dimensional solid or surface representation, thereby fully defining the external geometry of the object.\cite[p.~4]{gibson_additive_2010}

A fundamental distinction can be drawn between software that focuses primarily on surface representation and that which supports solid modelling.
For additive manufacturing, a fully enclosed model, \ie one that encloses a complete volume, is essential.
This requirement is often difficult to fulfil when using surface-modelling software, as surfaces are described by boundary curves that can be challenging to join with the precision and continuity necessary to form a closed body.
In contrast, software developed for manufacturing-oriented modelling is designed to maintain geometric and topological consistency, typically ensuring that resulting models are fully enclosed and possess a defined volume.\cite[p.~21]{gibson_additive_2010}
Although such software frequently includes additional functionalities, such as computer-aided manufacturing and physical simulations, only the computer-aided design aspect is relevant to the workflow considered in this study.

Solid-modelling systems used in manufacturing contexts are based on precise geometric definitions.
Typically, the process begins with two-dimensional sketches that incorporate parametric constraints specifying dimensions, angles, and positional relationships among components.
These sketches can be positioned on any plane within the coordinate system or on the surface of existing geometry.
Subsequent operations, such as extrusion, cutting, or lofting, then transform these sketches into three-dimensional volumes.

Modern CAD systems commonly rely on boundary representations that utilise Non-Uniform Rational B-Splines (NURBS) to mathematically define surfaces and curves.\cite[p.~21]{gibson_additive_2010}
To facilitate downstream applications such as slicing, these systems also provide tools for exporting models into a wide range of common interchange file formats.


\section{Slicer}

The fundamental principle of 3D printing is that components are created through the additive process of layering material, with each layer representing a thin cross-section of the part derived from the original CAD data.\cite[p.~2]{gibson_additive_2010}
The task of dividing a model into these discrete layers and generating sequential printing instructions corresponds to the second and third steps of the aforementioned workflow, and constitutes the primary function of the software tool commonly referred to as a slicer.
Although slicer applications may include limited features for creating or editing 3D models, their focus lies in preparing a completed model for printing and translating the three-dimensional data into printer-specific instructions.\cite{hudson_understanding_2016}

Slicers are capable of importing 3D models in a range of file formats and transforming them to the software's internal representation. 
Within the slicer interface, the model's individual components are placed on a virtual depiction of the printer's build platform or \textit{print bed}, the physical surface on which the object will be produced. 
Since printers have varying build volumes, the print bed measurements as well as the available height need to be available for correct part placement and to optimise print efficiency and material usage. 
While users may configure these parameters manually, most slicers provide a predefined library of printer profiles, often including models from multiple manufacturers, to facilitate setup.

In extrusion-based 3D printing, material characteristics such as rheology (the flow behaviour of molten materials) or melting temperature are as relevant as machine configurations, including print speed, nozzle geometry, and heating profiles.\cite{subbaraman_shape_2025}
These variables introduce a considerable degree of complexity to the printing process, as each parameter can influence the final print quality in different ways. 
For the purpose of conceptual clarity, this paper will differentiate between \textit{filament} and\textit{ printing parameters}, even though this distinction is not necessarily reflected in the slicer software's user interface. 

\subsection{Filament Presets}

3D printers support a wide variety of filament types, each with distinct material properties and potentially manufacturer-specific characteristics.
As a result, the associated printing parameters can vary considerably between materials.
To reduce this complexity, slicer software typically provides a selection of predefined filament presets, which may be either generic or tailored to specific manufacturers.
A preset is a collection of parameters that can be applied collectively to streamline the configuration process, ensuring that all model components, that the corresponding filament is assigned to, are treated uniformly.
These parameters generally include the filament’s name and base material, as well as recommended settings such as nozzle temperatures, extrusion speed, and cooling behaviour.
Users may duplicate and modify these presets to fine-tune individual parameters according to their specific printer setup or desired print quality, thereby enabling the reuse of optimised configurations across multiple projects.

\subsection{Printing Parameters}

Additional parameters frequently adjusted during print preparation include those related to print quality (\eg layer height, line width, seam positioning, precision, and bridging behaviour) as well as parameters that influence the model’s mechanical strength (\eg number of wall loops, shell thickness, and infill density or pattern).
Further settings determine the type and placement of support structures, including threshold angles at which supports are generated.
Unlike filament presets, these parameters are not necessarily tied to the selected material and may vary between individual model components.
Slicer software typically provides a number of standard presets for such configurations as well, while also allowing users to modify or create custom profiles as needed.

\subsection{Modifiers}

As previously stated in the introductory example, specific components within a 3D model may act as a modifier. 
They typically change printing parameters, such as the number of wall loops or infill density, locally for a limited region of the model.
For example, additional reinforcements may be applied to the area around the snap-fit closure, which is subjected to repeated mechanical stress, by increasing the number of wall loops.

Although it is possible to add such objects within the slicer's user interface, this approach tends to be cumbersome and imprecise.
Therefore users may prefer to define modifying objects during the 3D modelling stage, thus avoiding repetitive manual positioning in the slicer during each design iteration.
It is important to note that the term `modifier' is used this paper as an umbrella term encompassing not only conventional modifiers but also support blockers and enforcers.

\subsection{Slicer Files}

Once the model is sufficiently configured, the slicer converts it into a series of discrete layers, described as paths for the printer to follow. 
Due to the nature of this process, slicers typically manage two distinct file types. 
Before slicing, the 3D model is stored as a project file including geometric data as well as the associated parameters, allowing for the configuration to be reopened, modified and reused at a later point, independent of a specific printer. 
For this purpose, slicers generally use either the open-source format 3MF or a proprietary file format. 
Following the slicing process, the resulting model is exported in a format specific to a single printer ecosystem. 
This format is commonly referred to as G-code.\cite{kumar_printing_2023}

A variety of slicers are available nowadays, both general ones and versions customised by printer manufacturers. 
Since getting all the settings right for a particular 3D printer can be complex, some manufacturers prefer to adapt one of the open-source slicers, creating a version for their machines.\cite[p.~58]{horvath_mastering_2018} 
Thus, the majority of commonly used slicers share the same code base, derived from an ancestor called \textit{Slic3r}. 
The common lineage and the need to keep these projects open-source\footnote{\textit{Slic3r} is licensed under the \textit{GNU Affero General Public License}, version 3, which requires derivative works to remain open-source.}, has created a collaborative environment in which new features introduced in one version are oftentimes ported to others as well. 

\section{File Formats used in 3D Printing Context}

A plethora of alternative file formats are available to represent 3D models in general, however not all of these are equally appropriate for digital fabrication and only a select few have been adopted by the AM community. 
Hence, the following paragraphs will delve into the three most prevalent formats utilised in 3D printing.

\subsection{STL}
\label{chap:stl}

The STL (``STereoLithography'') format was introduced in the 1980s and has since gained widespread support among software applications used in digital fabrication. 
It is considered to be the most prominent and prevalent mesh file format, despite its stringent limitations.
The contents can be either encoded in binary or in ASCII text format. 
The latter is less common, but it is used for illustration and educational purposes. 
It describe an approximation of the surface geometry in the form of a mesh of triangles, each identified by a unit normal vector as well as three vertices, as seen in \ref{lst:stl_example}.
The normal vector is equal to the length ``1'', however it does not specify whether the dimensions are millimetres, inches or another unit entirely. 
Consequently, while the proportions of a model remain constant, the overall size depends on the user's knowledge about the model. 
Due to its simplicity, the file does not contain information regarding color, text or other common CAD model attributes.\cite[p.~341 ff.]{gibson_additive_2010}\cite{kumar_printing_2023}

Furthermore, since the specification requires each triangle to share two of its vertices with the triangle adjacent to it, thereby creating one continuous solid model, technically there will be exactly one body in the resulting file. 
For a model comprised of multiple parts, the user therefore is required to choose whether to combine all bodies into one or export each body as a separate file.
However, modern slicer software is capable of separating parts into disparate bodies, even if the user chooses to export a multi-body model into one file. 


\subsection{OBJ}
\label{chap:obj}

The OBJ file format is considered to be the second most common file format as it contains color and texture information in addition to the options STL presents.
It was developed by Wavefront Technologies and later introduced to various other 3D graphics systems due to its simplicity and open-source license.
Like STL files, it supports both ASCII text as well as binary encoding. 
While the 3D model does approximate the surface of the objects similar to STL files, OBJ also supports free-form surfaces and curves.\cite{kumar_printing_2023}

As can be seen in \ref{lst:obj_3d_model}, the model is represented by vertices as well. 
Each group of vertices is prefaced by a body name and optionally a material, referencing a separate \texttt{.mtl} (material template library) file. 

Common problems with this format include the file's complexity, making it difficult to repair non-closed surfaces, as well as the need for additional plug-ins in some CAD systems to export it. 
On top of that, having more than one file to describe a single 3D model is a potential cause of issues, especially when sharing, importing and transferring files. 


\subsection{3MF}
\label{chap:3mf}

3MF (``3D manufacturing format'') has been developed by the 3MF Consortium, founded by Microsoft, Autodesk, HP, 3D Systems, Prusa Research and several other key players in the field, to create a 3D printing file format that addresses the shortcomings of previous standards.
In particular the possibility for extensions and the support of colors as well as several different 3D model representations are aspects in which the format is specifically designed to support interoperability between CAD and AM systems.\cite{kumar_printing_2023}

\begin{figure}[htb]
    \dirtree{%
        .1 example.3mf.
        .2 \detokenize{_rels}/.
        .2 \detokenize{[Content_Types].xml}/.
        .2 3D/.
        .3 3dmodel.model.
        .2 Metadata/.
        .3 thumbnail.png.
    }
    \caption{Typical internal structure of 3MF file as exported from \emph{Autodesk Fusion}}
    \label{doc_structure}
\end{figure}

The two formats described above can generally be opened by a standard text editor to get a human readable result. 
3MF on the other hand is ``a ZIP archive that contains a complex file structure, defining a 3D model in Extensible Markup Language (XML).''\cite{rossel_security_2023}
Same as the other two formats, the 3D data is described by vertices, and triangles referencing three vertices each.
In addition, the XML file contains \texttt{colorgroup} tags in its resources section, that represent materials as RGBA color codes in hexadecimal format. 
Multiple 3D model files can be included, one of which needs to be identified in the \texttt{StartPart} relationship as the root model.

The \texttt{object} tags each contain attributes (body id, type, color id) as well as one triangular geometry mesh representation of an object volume. 
The latter is represented by vertices (x-, y-, z-coordinates in the unit specified within the model's attributes) and triangles referencing three vertices each oriented with normals pointing towards the exterior of their object.\cite{3mf_consortium_3mf_2025}
Meshes of different objects may intersect (as seen in the above mentioned example's case of wall loop modifiers) or self-intersect, in which case it must be specified which areas are inside or respectively outside the volume. 

\section{Software-specific Workflow}

The overall workflow doesn't vary greatly depending on specific software, however a few details should be noted. 
To explore the concept presented in this paper, \textit{Autodesk Fusion} was used as the 3D modelling software and \textit{OrcaSlicer} as slicer software.

\subsection{\textit{Autodesk Fusion}}

\textit{Autodesk Fusion}, released in 2013, offers various functionalities for computer-aided engineering (CAE) as well as product data management. 
The software's CAE environment includes tools for designing, generative design, rendering, and simulation, thus supporting the entire product development process with support for a variety of digital fabrication techniques\cite[p.~2]{junk_autodesk_2025}.  

Of particular relevance to this study are Fusion’s options for assigning \textit{physical materials} and \textit{appearances} to model bodies.
Physical materials can be applied to entire bodies to define their thermal, mechanical, strength and other engineering properties, particularly in regards to simulations, while also determining their visual appearance. 
Appearances, by contrast, can be assigned either to a body or individual faces, and affect only the visual rendering, overriding color and texture of previously applied materials. 

Fusion projects are stored in the cloud by default, making it rare for users to directly interact with the proprietary \texttt{.f3d} file format. 
Generally, files are only handled locally, when exporting the model from the Fusion ecosystem.
Two user interfaces provide this functionality: a general export dialogue tailored for all kinds of digital fabrication, and a 3D Print dialogue, tailored for additive manufacturing in particular.
Each offers a different set of file formats corresponding to its intended purpose.
The 3D Print dialogue allows exports in STL, 3MF, and OBJ formats, with the additional option to export either a single file containing all selected bodies, or separate files for each body.
Furthermore, for STL and OBJ exports, the user can specify the unit type (\eg millimetres or inches), whereas in 3MF exports this selection is predefined and therefore disabled.

\subsection{\textit{OrcaSlicer}}

\textit{OrcaSlicer} is one of the open-source slicers, originally forked from \textit{Bambu Studio}. 
It inherits the core functionality of its predecessor while introducing community-driven improvements, since it is a purely community-based project with no involvement from a 3D printer manufacturer.

The workflow within \textit{OrcaSlicer} changes subtly depending on the file format used for import.
When importing a model in STL format, all bodies combined into a single file are treated as one object.
If the user intends to print different parts in separate colours or apply modifiers, such as strengthening regions around a snap-fit closure, an additional step is required to split the model into separate objects within the slicer.

\begin{wrapfigure}[17]{l}{0.45\linewidth}
    \vspace{-13pt}
  \includegraphics[width=\linewidth]{obj_color_dialog.png}
  \label{fig:obj_color_dialog}
  \caption{OBJ color mapping dialog in \textit{OrcaSlicer}}
\end{wrapfigure}

When importing multiple STL files at once, \textit{OrcaSlicer} prompts the user to decide whether the files should be loaded as a single object with multiple parts or as separate objects.
Loading them as a single object maintains their spatial relationships, which is useful when dealing with models that include modifiers or precisely aligned parts.
In contrast, individual objects are each positioned with one surface flush to the build plate, which can lead to misalignment if the original relative positioning must be preserved.

When importing an OBJ file, \textit{OrcaSlicer} displays a dedicated dialogue asking the user how to map the file’s embedded colour information to existing or new filament presets.
This feature enables users to align the software’s filament colours (i.e., the visual representations within \textit{OrcaSlicer}, not necessarily corresponding to the physical filament) with the colours assigned during the modelling stage in Fusion.

Importing a 3MF file containing multiple objects, \textit{OrcaSlicer} prompts user to decide whether to treat the file as a single object with multiple parts or as several independent objects, comparable to the behaviour when importing multiple STL or OBJ files at once.
However, unlike the OBJ import, \textit{OrcaSlicer} does not offer the user an option to map RGBA colour data embedded in the 3MF file to existing or new filament presets.
Consequently, all imported parts are initially assigned the primary filament colour by default, requiring users to manually reapply materials.
This limitation has been criticised by users, for example in \textit{OrcaSlicer issue \#2438}\footnote{Single Object Multiple parts with Multiple Colours: [...] I have made a number of 3D models with multiple parts in single 3mf file. I mostly use Fusion 360 for models and during modelling I assign colours to each part. But when I upload this 3mf file to the slicer, It assigns the primary filament colour to all the parts. I was wondering if the slicer would be able to distinguish the parts with different colours and assign either the original colour to the parts or even randon colours(which can be changed later).
For example if my 3mf file has 5 colour parts, the slicer should be to detect them. Then it should find which colour has the maximum spread by area and assign the filament 1 to it. The second largest should be assigned the second filament and so on till the last.[...] \\ \url{https://github.com/SoftFever/OrcaSlicer/discussions/2438}}.

Within \textit{OrcaSlicer}, materials are assigned using system-defined or user-defined filament presets.
The software proactively prompts users to save any modifications made to material presets, allowing for user-friendly management of customised presets tailored either to specific printers, filament types, or even individual projects.
By default, when a new model or project is imported, OrcaSlicer automatically applies the first available user-defined filament preset.

Printing presets are handled similarly, however is has become apparent that these can only be applied globally on project level, rather than to specific individual objects within a project. 
This limitation becomes particularly noticeable in complex assemblies comprised of multiple components requiring different settings,
as highlighted in \textit{OrcaSlicer Issue \#7769}\footnote{Copy/paste per-object settings: When planning/building large projects in-slicer it often becomes an extremely tedious task to click all the objects you want a specific set of settings to apply to. The larger the assembly/project you are planning, the easier it is to miss a body and make a mistake. If you "assemble" models together you cannot edit multiple sub-objects across assemblies.\\ \url{https://github.com/SoftFever/OrcaSlicer/issues/7769}}.
As a result, users often resort to repetitive manual adjustments when handling multi-part models with diverse printing requirements.

% \section{Related Work and Available Solutions}

% ``CAM-based design tools have emerged where toolpaths are specified directly. This approach is contrasted with CAD-based tools which are based on geometric models [23]. CAM-based tools naturally shift focus to the machine and its settings. Prior work generates custom toolpaths using traditional CAD environments [20 , 21 ], code [ 25 , 54 , 62 ], hand-drawn sketches [23], and domain-specific abstractions [ 5 ].''[its not the shape]

% how do others solve the issue $\rightarrow$ focus on documentation [its not the shape]

% $\rightarrow$ maybe also refer to git issues here 